\documentclass[12pt,a4paper]{article}
\usepackage[utf8]{vietnam}
\usepackage{amsmath,amssymb}
\usepackage{enumerate}
\usepackage{tabularx}
\usepackage{tikz}
\usepackage{xcolor}
\usepackage[top=0.45cm, bottom=0.45cm, left=0.8cm, right=0.8cm]{geometry}
\DeclareUnicodeCharacter{00B2}{\ensuremath{^2}}
\DeclareUnicodeCharacter{00B3}{\ensuremath{^3}}

\setlength{\parskip}{0pt}
\setlength{\parindent}{0pt}
\setlength{\fboxsep}{3.5pt}
\setlength{\fboxrule}{0.5pt}

% Dòng kẻ chấm mờ cự ly chuẩn viết tay tiểu học (18pt)
\newcommand{\dongke}[1]{%
  \begingroup%
  \setlength{\parskip}{0pt}%
  \baselineskip=18pt%
  \vspace{0.4mm}%
  \foreach \i in {1,...,#1}{%
    \noindent{\color{gray!55}\dotfill}\par%
  }%
  \vspace{0.2mm}%
  \endgroup%
}

% Khung bao quanh bài tập tự luận
\newcommand{\khungBaiTap}[1]{%
  \vspace{0.8mm}%
  \noindent\fbox{%
    \parbox{\dimexpr\linewidth-2\fboxsep-2\fboxrule\relax}{%
      \vspace{0.3mm}%
      #1%
      \vspace{0.3mm}%
    }%
  }%
  \vspace{0.5mm}%
}

% Ô bài tập đặt tính
\newcommand{\oDatTinh}[3][4]{%
  \vspace{0.3mm}%
  \noindent\textbf{#2} #3\par%
  \dongke{#1}%
}

% Câu trắc nghiệm dàn 4 đáp án trên 1 hàng ngang
\newcommand{\cauTN}[6]{%
  \noindent\textbf{Câu #1.} #2\par\vspace{0.1mm}%
  \noindent%
  \makebox[0.25\linewidth][l]{\textbf{A.} #3}%
  \makebox[0.25\linewidth][l]{\textbf{B.} #4}%
  \makebox[0.25\linewidth][l]{\textbf{C.} #5}%
  \makebox[0.25\linewidth][l]{\textbf{D.} #6}\par\vspace{0.6mm}%
}

% Câu trắc nghiệm dàn 2 đáp án / hàng
\newcommand{\cauTNhaiDong}[6]{%
  \noindent\textbf{Câu #1.} #2\par\vspace{0.1mm}%
  \noindent%
  \makebox[0.5\linewidth][l]{\textbf{A.} #3}%
  \makebox[0.5\linewidth][l]{\textbf{B.} #4}\par\vspace{0.1mm}%
  \noindent%
  \makebox[0.5\linewidth][l]{\textbf{C.} #5}%
  \makebox[0.5\linewidth][l]{\textbf{D.} #6}\par\vspace{0.6mm}%
}

\begin{document}
\pagestyle{empty}
\fontsize{10.2pt}{13.2pt}\selectfont

% =============================================================================
% ========================= MÃ ĐỀ: 401 ========================================
% =============================================================================

% -------------------- TRANG 1: MẶT TRƯỚC --------------------
\noindent
\begin{tabularx}{\linewidth}{|X|c|}
\hline
\textbf{PHIẾU BÀI TẬP CUỐI TUẦN -- TOÁN LỚP 4 (TUẦN 09)} & \textbf{MÃ ĐỀ: 401} \\
\textit{Trung tâm Toán Thầy Lĩnh -- Năm học: 2026 - 2027} & \textit{Thời gian: 45 phút} \\
\hline
\end{tabularx}

\vspace{1mm}

\noindent
\begin{tabularx}{\linewidth}{|X|p{3.5cm}|}
\hline
\rule[-2mm]{0pt}{7.5mm}\textbf{Họ và tên:} \dotfill & \textbf{Lớp:} \makebox[2cm]{\dotfill} \\
\hline
\end{tabularx}

\vspace{1mm}

\noindent
\begin{tabularx}{\linewidth}{|p{3cm}|X|}
\hline
\centering\textbf{ĐIỂM SỐ} & \centering\textbf{LỜI NHẬN XÉT CỦA THẦY CÔ} \tabularnewline
\hline
\rule{0pt}{1.1cm} & \tabularnewline
\hline
\end{tabularx}

\vspace{1mm}

\noindent\textbf{BẢNG ĐIỀN ĐÁP ÁN TRẮC NGHIỆM (Mỗi câu đúng 0,25 điểm):}
\vspace{0.6mm}

\begin{tabularx}{\linewidth}{|c|*{12}{>{\centering\arraybackslash}X|}}
\hline
\textbf{Câu} & \textbf{1} & \textbf{2} & \textbf{3} & \textbf{4} & \textbf{5} & \textbf{6} & \textbf{7} & \textbf{8} & \textbf{9} & \textbf{10} & \textbf{11} & \textbf{12} \\
\hline
\textbf{Đáp án} & & & & & & & & & & & & \\
\hline
\end{tabularx}

\vspace{1.2mm}

\noindent\textbf{\large PHẦN I. TRẮC NGHIỆM KHÁCH QUAN (3,0 điểm)}

\cauTN{1}{Số gồm $5$ chục nghìn, $8$ trăm và $3$ đơn vị viết là:}{$50\,803$}{$58\,003$}{$50\,830$}{$508\,003$}
\cauTN{2}{Chữ số $7$ trong số $67\,425$ thuộc hàng nào?}{Hàng chục}{Hàng trăm}{Hàng nghìn}{Hàng chục nghìn}
\cauTN{3}{Số liền trước của số $45\,000$ là:}{$44\,999$}{$45\,001$}{$44\,990$}{$44\,900$}
\cauTN{4}{Số lớn nhất trong các số $78\,542;\; 78\,524;\; 78\,254;\; 78\,452$ là:}{$78\,524$}{$78\,542$}{$78\,452$}{$78\,254$}
\cauTN{5}{Làm tròn số $43\,682$ đến hàng trăm ta được số:}{$43\,600$}{$43\,700$}{$44\,000$}{$43\,680$}
\cauTN{6}{Trong các số sau, số nào là số chẵn?}{$25\,431$}{$38\,705$}{$46\,890$}{$52\,117$}
\cauTN{7}{Giá trị của biểu thức $m + n$ với $m = 25\,000$ và $n = 35\,000$ là:}{$50\,000$}{$60\,000$}{$70\,000$}{$55\,000$}
\cauTN{8}{Số thích hợp điền vào chỗ chấm $4\,\text{m}\,6\,\text{cm} = \dots\,\text{cm}$ là:}{$46$}{$406$}{$460$}{$4\,006$}
\cauTN{9}{Khối lượng $5\,\text{kg}\,300\,\text{g} = \dots\,\text{g}$ là:}{$530\,\text{g}$}{$5\,030\,\text{g}$}{$5\,300\,\text{g}$}{$50\,300\,\text{g}$}
\cauTN{10}{Một năm nhuận có bao nhiêu ngày?}{$364$ ngày}{$365$ ngày}{$366$ ngày}{$367$ ngày}
\cauTNhaiDong{11}{Một hình vuông có chu vi là $36\,\text{cm}$. Diện tích của hình vuông đó là:}{$81\,\text{cm}^2$}{$36\,\text{cm}^2$}{$72\,\text{cm}^2$}{$18\,\text{cm}^2$}
\cauTNhaiDong{12}{Mua $4$ quyển vở cùng loại hết $28\,000$ đồng. Mua $6$ quyển vở như thế hết số tiền là:}{$35\,000$ đồng}{$42\,000$ đồng}{$48\,000$ đồng}{$56\,000$ đồng}

\vspace{0.4mm}
\noindent\textbf{\large PHẦN II. TỰ LUẬN (7,0 điểm)}

\noindent\textbf{Bài 1 (2,0 điểm). Đặt tính rồi tính:}
\vspace{0.4mm}

\noindent
\begin{tabularx}{\linewidth}{|X|X|X|X|}
\hline
\oDatTinh[4]{a)}{$34\,528 + 28\,345$} & 
\oDatTinh[4]{b)}{$72\,615 - 38\,254$} & 
\oDatTinh[4]{c)}{$14\,216 \times 4$} & 
\oDatTinh[4]{d)}{$48\,725 : 5$} \tabularnewline
\hline
\end{tabularx}

\newpage
% -------------------- TRANG 2: MẶT SAU --------------------
\noindent
\begin{tabularx}{\linewidth}{|X|c|}
\hline
\textbf{PHẦN TỰ LUẬN (TIẾP THEO) -- PHIẾU BÀI TẬP TUẦN 09} & \textbf{MÃ ĐỀ: 401} \\
\hline
\end{tabularx}

\vspace{0.8mm}

% BÀI 2: TÌM X (4 Ý)
\khungBaiTap{%
  \textbf{Bài 2 (1,5 điểm). Tìm $x$:}\par\vspace{0.6mm}
  \noindent
  \begin{minipage}[t]{0.485\linewidth}
    \textbf{a)} $x - 14\,250 = 32\,480$\par
    \dongke{3}
    \vspace{0.4mm}
    \textbf{c)} $x \times 6 = 48\,624$\par
    \dongke{3}
  \end{minipage}\hfill
  \begin{minipage}[t]{0.485\linewidth}
    \textbf{b)} $58\,420 - x = 24\,160$\par
    \dongke{3}
    \vspace{0.4mm}
    \textbf{d)} $x : 7 = 8\,142$\par
    \dongke{3}
  \end{minipage}
}

\vspace{0.8mm}

% BÀI 3: ĐỔI ĐƠN VỊ (4 Ý - 1 CỘT TỪ TRÊN XUỐNG)
\khungBaiTap{%
  \textbf{Bài 3 (1,0 điểm). Viết số thích hợp vào chỗ chấm:}\par\vspace{0.8mm}
  \noindent
  a) $7\,\text{km}\,25\,\text{m} = \dotfill\,\text{m}$\par\vspace{1.4mm}
  \noindent
  b) $6\,\text{kg}\,450\,\text{g} = \dotfill\,\text{g}$\par\vspace{1.4mm}
  \noindent
  c) $8\,\text{m}^2\,15\,\text{dm}^2 = \dotfill\,\text{dm}^2$\par\vspace{1.4mm}
  \noindent
  d) $3\,\text{phút}\,40\,\text{giây} = \dotfill\,\text{giây}$\par
}

\vspace{0.8mm}

% BÀI 4: BÀI TOÁN GIẢI (11 DÒNG KẺ)
\khungBaiTap{%
  \textbf{Bài 4 (1,5 điểm).} Một thửa ruộng hình chữ nhật có chiều dài $120\,m$, chiều rộng $80\,m$. Người ta cấy lúa trên thửa ruộng đó, cứ $100\,m^2$ thu hoạch được $65\,kg$ thóc. Hỏi cả thửa ruộng thu hoạch được bao nhiêu tạ thóc? Diện tích thửa ruộng bằng bao nhiêu héc-ta?
  \par\vspace{0.3mm}%
  \noindent\textit{Bài giải:}
  \dongke{11}
}

\vspace{0.8mm}

% BÀI 5: TOÁN NÂNG CAO (8 DÒNG KẺ)
\khungBaiTap{%
  \textbf{Bài 5 (1,0 điểm).} Tính bằng cách thuận tiện nhất:\par\vspace{0.6mm}\centerline{$K = 0{,}25 \times 7{,}4 \times 4 + 12{,}6 \times 5$}
  \par\vspace{0.3mm}%
  \noindent\textit{Bài làm:}
  \dongke{8}
}

\newpage


% =============================================================================
% ========================= MÃ ĐỀ: 402 ========================================
% =============================================================================

% -------------------- TRANG 1: MẶT TRƯỚC --------------------
\noindent
\begin{tabularx}{\linewidth}{|X|c|}
\hline
\textbf{PHIẾU BÀI TẬP CUỐI TUẦN -- TOÁN LỚP 4 (TUẦN 09)} & \textbf{MÃ ĐỀ: 402} \\
\textit{Trung tâm Toán Thầy Lĩnh -- Năm học: 2026 - 2027} & \textit{Thời gian: 45 phút} \\
\hline
\end{tabularx}

\vspace{1mm}

\noindent
\begin{tabularx}{\linewidth}{|X|p{3.5cm}|}
\hline
\rule[-2mm]{0pt}{7.5mm}\textbf{Họ và tên:} \dotfill & \textbf{Lớp:} \makebox[2cm]{\dotfill} \\
\hline
\end{tabularx}

\vspace{1mm}

\noindent
\begin{tabularx}{\linewidth}{|p{3cm}|X|}
\hline
\centering\textbf{ĐIỂM SỐ} & \centering\textbf{LỜI NHẬN XÉT CỦA THẦY CÔ} \tabularnewline
\hline
\rule{0pt}{1.1cm} & \tabularnewline
\hline
\end{tabularx}

\vspace{1mm}

\noindent\textbf{BẢNG ĐIỀN ĐÁP ÁN TRẮC NGHIỆM (Mỗi câu đúng 0,25 điểm):}
\vspace{0.6mm}

\begin{tabularx}{\linewidth}{|c|*{12}{>{\centering\arraybackslash}X|}}
\hline
\textbf{Câu} & \textbf{1} & \textbf{2} & \textbf{3} & \textbf{4} & \textbf{5} & \textbf{6} & \textbf{7} & \textbf{8} & \textbf{9} & \textbf{10} & \textbf{11} & \textbf{12} \\
\hline
\textbf{Đáp án} & & & & & & & & & & & & \\
\hline
\end{tabularx}

\vspace{1.2mm}

\noindent\textbf{\large PHẦN I. TRẮC NGHIỆM KHÁCH QUAN (3,0 điểm)}

\cauTN{1}{Số gồm $7$ chục nghìn, $4$ trăm và $6$ đơn vị viết là:}{$70\,406$}{$74\,006$}{$70\,460$}{$704\,006$}
\cauTN{2}{Chữ số $8$ trong số $84\,315$ thuộc hàng nào?}{Hàng chục}{Hàng trăm}{Hàng nghìn}{Hàng chục nghìn}
\cauTN{3}{Số liền sau của số $59\,999$ là:}{$60\,000$}{$59\,998$}{$60\,001$}{$59\,900$}
\cauTN{4}{Số bé nhất trong các số $62\,845;\; 62\,485;\; 62\,584;\; 62\,854$ là:}{$62\,845$}{$62\,485$}{$62\,584$}{$62\,854$}
\cauTN{5}{Làm tròn số $56\,248$ đến hàng trăm ta được số:}{$56\,200$}{$56\,300$}{$56\,000$}{$56\,250$}
\cauTN{6}{Trong các số sau, số nào là số lẻ?}{$42\,680$}{$37\,845$}{$56\,922$}{$78\,134$}
\cauTN{7}{Giá trị của biểu thức $a - b$ với $a = 80\,000$ và $b = 25\,000$ là:}{$45\,000$}{$55\,000$}{$65\,000$}{$50\,000$}
\cauTN{8}{Số thích hợp điền vào chỗ chấm $6\,\text{m}\,8\,\text{cm} = \dots\,\text{cm}$ là:}{$68$}{$608$}{$680$}{$6\,008$}
\cauTN{9}{Khối lượng $4\,\text{kg}\,500\,\text{g} = \dots\,\text{g}$ là:}{$450\,\text{g}$}{$4\,050\,\text{g}$}{$4\,500\,\text{g}$}{$45\,000\,\text{g}$}
\cauTN{10}{Thế kỉ XXI bắt đầu từ năm nào đến năm nào?}{Từ năm 2000 đến 2099}{Từ năm 2001 đến 2100}{Từ năm 1901 đến 2000}{Từ năm 2000 đến 2100}
\cauTNhaiDong{11}{Một hình chữ nhật có chiều dài $12\,\text{cm}$, chiều rộng $8\,\text{cm}$. Diện tích là:}{$40\,\text{cm}^2$}{$96\,\text{cm}^2$}{$20\,\text{cm}^2$}{$48\,\text{cm}^2$}
\cauTNhaiDong{12}{Mua $5$ chiếc bút cùng loại hết $35\,000$ đồng. Mua $8$ chiếc bút như thế hết số tiền là:}{$49\,000$ đồng}{$56\,000$ đồng}{$63\,000$ đồng}{$42\,000$ đồng}

\vspace{0.4mm}
\noindent\textbf{\large PHẦN II. TỰ LUẬN (7,0 điểm)}

\noindent\textbf{Bài 1 (2,0 điểm). Đặt tính rồi tính:}
\vspace{0.4mm}

\noindent
\begin{tabularx}{\linewidth}{|X|X|X|X|}
\hline
\oDatTinh[4]{a)}{$42\,637 + 35\,248$} & 
\oDatTinh[4]{b)}{$84\,520 - 46\,275$} & 
\oDatTinh[4]{c)}{$16\,315 \times 5$} & 
\oDatTinh[4]{d)}{$36\,848 : 4$} \tabularnewline
\hline
\end{tabularx}

\newpage
% -------------------- TRANG 2: MẶT SAU --------------------
\noindent
\begin{tabularx}{\linewidth}{|X|c|}
\hline
\textbf{PHẦN TỰ LUẬN (TIẾP THEO) -- PHIẾU BÀI TẬP TUẦN 09} & \textbf{MÃ ĐỀ: 402} \\
\hline
\end{tabularx}

\vspace{0.8mm}

% BÀI 2: TÌM X (4 Ý)
\khungBaiTap{%
  \textbf{Bài 2 (1,5 điểm). Tìm $x$:}\par\vspace{0.6mm}
  \noindent
  \begin{minipage}[t]{0.485\linewidth}
    \textbf{a)} $x - 18\,360 = 25\,420$\par
    \dongke{3}
    \vspace{0.4mm}
    \textbf{c)} $x \times 5 = 45\,325$\par
    \dongke{3}
  \end{minipage}\hfill
  \begin{minipage}[t]{0.485\linewidth}
    \textbf{b)} $67\,850 - x = 32\,410$\par
    \dongke{3}
    \vspace{0.4mm}
    \textbf{d)} $x : 6 = 7\,215$\par
    \dongke{3}
  \end{minipage}
}

\vspace{0.8mm}

% BÀI 3: ĐỔI ĐƠN VỊ (4 Ý - 1 CỘT TỪ TRÊN XUỐNG)
\khungBaiTap{%
  \textbf{Bài 3 (1,0 điểm). Viết số thích hợp vào chỗ chấm:}\par\vspace{0.8mm}
  \noindent
  a) $5\,\text{km}\,40\,\text{m} = \dotfill\,\text{m}$\par\vspace{1.4mm}
  \noindent
  b) $8\,\text{kg}\,250\,\text{g} = \dotfill\,\text{g}$\par\vspace{1.4mm}
  \noindent
  c) $6\,\text{m}^2\,45\,\text{dm}^2 = \dotfill\,\text{dm}^2$\par\vspace{1.4mm}
  \noindent
  d) $4\,\text{phút}\,15\,\text{giây} = \dotfill\,\text{giây}$\par
}

\vspace{0.8mm}

% BÀI 4: BÀI TOÁN GIẢI (11 DÒNG KẺ)
\khungBaiTap{%
  \textbf{Bài 4 (1,5 điểm).} Một mảnh đất hình bình hành có độ dài đáy $150\,m$, chiều cao bằng $\frac{2}{3}$ độ dài đáy. Người ta trồng khoai trên mảnh đất đó, trung bình mỗi héc-ta thu được $12$ tấn khoai. Hỏi cả mảnh đất thu hoạch được bao nhiêu tấn khoai? Diện tích mảnh đất bằng bao nhiêu héc-ta?
  \par\vspace{0.3mm}%
  \noindent\textit{Bài giải:}
  \dongke{11}
}

\vspace{0.8mm}

% BÀI 5: TOÁN NÂNG CAO (8 DÒNG KẺ)
\khungBaiTap{%
  \textbf{Bài 5 (1,0 điểm).} Tính bằng cách thuận tiện nhất:\par\vspace{0.6mm}\centerline{$K = 2{,}5 \times 3{,}6 + 6{,}4 \times 2{,}5$}
  \par\vspace{0.3mm}%
  \noindent\textit{Bài làm:}
  \dongke{8}
}

\newpage


% =============================================================================
% ========================= MÃ ĐỀ: 403 ========================================
% =============================================================================

% -------------------- TRANG 1: MẶT TRƯỚC --------------------
\noindent
\begin{tabularx}{\linewidth}{|X|c|}
\hline
\textbf{PHIẾU BÀI TẬP CUỐI TUẦN -- TOÁN LỚP 4 (TUẦN 09)} & \textbf{MÃ ĐỀ: 403} \\
\textit{Trung tâm Toán Thầy Lĩnh -- Năm học: 2026 - 2027} & \textit{Thời gian: 45 phút} \\
\hline
\end{tabularx}

\vspace{1mm}

\noindent
\begin{tabularx}{\linewidth}{|X|p{3.5cm}|}
\hline
\rule[-2mm]{0pt}{7.5mm}\textbf{Họ và tên:} \dotfill & \textbf{Lớp:} \makebox[2cm]{\dotfill} \\
\hline
\end{tabularx}

\vspace{1mm}

\noindent
\begin{tabularx}{\linewidth}{|p{3cm}|X|}
\hline
\centering\textbf{ĐIỂM SỐ} & \centering\textbf{LỜI NHẬN XÉT CỦA THẦY CÔ} \tabularnewline
\hline
\rule{0pt}{1.1cm} & \tabularnewline
\hline
\end{tabularx}

\vspace{1mm}

\noindent\textbf{BẢNG ĐIỀN ĐÁP ÁN TRẮC NGHIỆM (Mỗi câu đúng 0,25 điểm):}
\vspace{0.6mm}

\begin{tabularx}{\linewidth}{|c|*{12}{>{\centering\arraybackslash}X|}}
\hline
\textbf{Câu} & \textbf{1} & \textbf{2} & \textbf{3} & \textbf{4} & \textbf{5} & \textbf{6} & \textbf{7} & \textbf{8} & \textbf{9} & \textbf{10} & \textbf{11} & \textbf{12} \\
\hline
\textbf{Đáp án} & & & & & & & & & & & & \\
\hline
\end{tabularx}

\vspace{1.2mm}

\noindent\textbf{\large PHẦN I. TRẮC NGHIỆM KHÁCH QUAN (3,0 điểm)}

\cauTN{1}{Số gồm $9$ chục nghìn, $5$ trăm và $8$ đơn vị viết là:}{$90\,508$}{$95\,008$}{$90\,580$}{$905\,008$}
\cauTN{2}{Chữ số $5$ trong số $45\,230$ thuộc hàng nào?}{Hàng trăm}{Hàng chục}{Hàng nghìn}{Hàng chục nghìn}
\cauTN{3}{Số liền trước của số $70\,000$ là:}{$69\,999$}{$70\,001$}{$69\,990$}{$69\,900$}
\cauTN{4}{Số lớn nhất trong các số $85\,412;\; 85\,421;\; 85\,142;\; 85\,214$ là:}{$85\,412$}{$85\,421$}{$85\,142$}{$85\,214$}
\cauTN{5}{Làm tròn số $64\,785$ đến hàng nghìn ta được số:}{$64\,000$}{$65\,000$}{$64\,800$}{$64\,700$}
\cauTN{6}{Trong các số sau, số nào là số chẵn?}{$51\,239$}{$63\,218$}{$47\,805$}{$39\,417$}
\cauTN{7}{Giá trị của biểu thức $x + y$ với $x = 36\,000$ và $y = 24\,000$ là:}{$50\,000$}{$60\,000$}{$70\,000$}{$55\,000$}
\cauTN{8}{Số thích hợp điền vào chỗ chấm $8\,\text{m}\,5\,\text{cm} = \dots\,\text{cm}$ là:}{$85$}{$805$}{$850$}{$8\,005$}
\cauTN{9}{Khối lượng $7\,\text{kg}\,600\,\text{g} = \dots\,\text{g}$ là:}{$760\,\text{g}$}{$7\,060\,\text{g}$}{$7\,600\,\text{g}$}{$76\,000\,\text{g}$}
\cauTN{10}{Một thế kỉ bằng bao nhiêu năm?}{$10$ năm}{$50$ năm}{$100$ năm}{$1\,000$ năm}
\cauTNhaiDong{11}{Một hình vuông có chu vi là $48\,\text{cm}$. Diện tích hình vuông đó là:}{$96\,\text{cm}^2$}{$144\,\text{cm}^2$}{$48\,\text{cm}^2$}{$24\,\text{cm}^2$}
\cauTNhaiDong{12}{Mua $6$ hộp sữa cùng loại hết $48\,000$ đồng. Mua $9$ hộp sữa như thế hết số tiền là:}{$64\,000$ đồng}{$72\,000$ đồng}{$80\,000$ đồng}{$56\,000$ đồng}

\vspace{0.4mm}
\noindent\textbf{\large PHẦN II. TỰ LUẬN (7,0 điểm)}

\noindent\textbf{Bài 1 (2,0 điểm). Đặt tính rồi tính:}
\vspace{0.4mm}

\noindent
\begin{tabularx}{\linewidth}{|X|X|X|X|}
\hline
\oDatTinh[4]{a)}{$53\,724 + 28\,169$} & 
\oDatTinh[4]{b)}{$91\,432 - 45\,286$} & 
\oDatTinh[4]{c)}{$15\,240 \times 6$} & 
\oDatTinh[4]{d)}{$56\,424 : 8$} \tabularnewline
\hline
\end{tabularx}

\newpage
% -------------------- TRANG 2: MẶT SAU --------------------
\noindent
\begin{tabularx}{\linewidth}{|X|c|}
\hline
\textbf{PHẦN TỰ LUẬN (TIẾP THEO) -- PHIẾU BÀI TẬP TUẦN 09} & \textbf{MÃ ĐỀ: 403} \\
\hline
\end{tabularx}

\vspace{0.8mm}

% BÀI 2: TÌM X (4 Ý)
\khungBaiTap{%
  \textbf{Bài 2 (1,5 điểm). Tìm $x$:}\par\vspace{0.6mm}
  \noindent
  \begin{minipage}[t]{0.485\linewidth}
    \textbf{a)} $x - 23\,450 = 41\,280$\par
    \dongke{3}
    \vspace{0.4mm}
    \textbf{c)} $x \times 7 = 56\,490$\par
    \dongke{3}
  \end{minipage}\hfill
  \begin{minipage}[t]{0.485\linewidth}
    \textbf{b)} $74\,650 - x = 28\,320$\par
    \dongke{3}
    \vspace{0.4mm}
    \textbf{d)} $x : 8 = 6\,312$\par
    \dongke{3}
  \end{minipage}
}

\vspace{0.8mm}

% BÀI 3: ĐỔI ĐƠN VỊ (4 Ý - 1 CỘT TỪ TRÊN XUỐNG)
\khungBaiTap{%
  \textbf{Bài 3 (1,0 điểm). Viết số thích hợp vào chỗ chấm:}\par\vspace{0.8mm}
  \noindent
  a) $9\,\text{km}\,60\,\text{m} = \dotfill\,\text{m}$\par\vspace{1.4mm}
  \noindent
  b) $5\,\text{kg}\,750\,\text{g} = \dotfill\,\text{g}$\par\vspace{1.4mm}
  \noindent
  c) $7\,\text{m}^2\,30\,\text{dm}^2 = \dotfill\,\text{dm}^2$\par\vspace{1.4mm}
  \noindent
  d) $2\,\text{phút}\,45\,\text{giây} = \dotfill\,\text{giây}$\par
}

\vspace{0.8mm}

% BÀI 4: BÀI TOÁN GIẢI (11 DÒNG KẺ)
\khungBaiTap{%
  \textbf{Bài 4 (1,5 điểm).} Một khu vườn hình chữ nhật có chu vi $500\,m$, chiều rộng $90\,m$. Người ta trồng cam trên khu vườn đó, cứ $100\,m^2$ thu được $80\,kg$ cam. Hỏi khu vườn thu hoạch được bao nhiêu tấn cam? Diện tích khu vườn bằng bao nhiêu héc-ta?
  \par\vspace{0.3mm}%
  \noindent\textit{Bài giải:}
  \dongke{11}
}

\vspace{0.8mm}

% BÀI 5: TOÁN NÂNG CAO (8 DÒNG KẺ)
\khungBaiTap{%
  \textbf{Bài 5 (1,0 điểm).} Tính bằng cách thuận tiện nhất:\par\vspace{0.6mm}\centerline{$K = 1{,}25 \times 4 + 2{,}5 \times 8 + 3{,}75 \times 4$}
  \par\vspace{0.3mm}%
  \noindent\textit{Bài làm:}
  \dongke{8}
}

\newpage


% =============================================================================
% ========================= MÃ ĐỀ: 404 ========================================
% =============================================================================

% -------------------- TRANG 1: MẶT TRƯỚC --------------------
\noindent
\begin{tabularx}{\linewidth}{|X|c|}
\hline
\textbf{PHIẾU BÀI TẬP CUỐI TUẦN -- TOÁN LỚP 4 (TUẦN 09)} & \textbf{MÃ ĐỀ: 404} \\
\textit{Trung tâm Toán Thầy Lĩnh -- Năm học: 2026 - 2027} & \textit{Thời gian: 45 phút} \\
\hline
\end{tabularx}

\vspace{1mm}

\noindent
\begin{tabularx}{\linewidth}{|X|p{3.5cm}|}
\hline
\rule[-2mm]{0pt}{7.5mm}\textbf{Họ và tên:} \dotfill & \textbf{Lớp:} \makebox[2cm]{\dotfill} \\
\hline
\end{tabularx}

\vspace{1mm}

\noindent
\begin{tabularx}{\linewidth}{|p{3cm}|X|}
\hline
\centering\textbf{ĐIỂM SỐ} & \centering\textbf{LỜI NHẬN XÉT CỦA THẦY CÔ} \tabularnewline
\hline
\rule{0pt}{1.1cm} & \tabularnewline
\hline
\end{tabularx}

\vspace{1mm}

\noindent\textbf{BẢNG ĐIỀN ĐÁP ÁN TRẮC NGHIỆM (Mỗi câu đúng 0,25 điểm):}
\vspace{0.6mm}

\begin{tabularx}{\linewidth}{|c|*{12}{>{\centering\arraybackslash}X|}}
\hline
\textbf{Câu} & \textbf{1} & \textbf{2} & \textbf{3} & \textbf{4} & \textbf{5} & \textbf{6} & \textbf{7} & \textbf{8} & \textbf{9} & \textbf{10} & \textbf{11} & \textbf{12} \\
\hline
\textbf{Đáp án} & & & & & & & & & & & & \\
\hline
\end{tabularx}

\vspace{1.2mm}

\noindent\textbf{\large PHẦN I. TRẮC NGHIỆM KHÁCH QUAN (3,0 điểm)}

\cauTN{1}{Số gồm $6$ chục nghìn, $2$ trăm và $9$ đơn vị viết là:}{$60\,209$}{$62\,009$}{$60\,290$}{$602\,009$}
\cauTN{2}{Chữ số $4$ trong số $34\,712$ thuộc hàng nào?}{Hàng trăm}{Hàng chục}{Hàng nghìn}{Hàng chục nghìn}
\cauTN{3}{Số liền sau của số $89\,999$ là:}{$90\,000$}{$89\,998$}{$90\,001$}{$89\,900$}
\cauTN{4}{Số bé nhất trong các số $43\,512;\; 43\,152;\; 43\,521;\; 43\,215$ là:}{$43\,512$}{$43\,152$}{$43\,521$}{$43\,215$}
\cauTN{5}{Làm tròn số $72\,384$ đến hàng trăm ta được số:}{$72\,300$}{$72\,400$}{$72\,000$}{$72\,380$}
\cauTN{6}{Trong các số sau, số nào là số lẻ?}{$62\,418$}{$74\,930$}{$54\,827$}{$88\,104$}
\cauTN{7}{Giá trị của biểu thức $m - n$ với $m = 90\,000$ và $n = 35\,000$ là:}{$45\,000$}{$55\,000$}{$65\,000$}{$50\,000$}
\cauTN{8}{Số thích hợp điền vào chỗ chấm $5\,\text{m}\,9\,\text{cm} = \dots\,\text{cm}$ là:}{$59$}{$509$}{$590$}{$5\,009$}
\cauTN{9}{Khối lượng $6\,\text{kg}\,800\,\text{g} = \dots\,\text{g}$ là:}{$680\,\text{g}$}{$6\,080\,\text{g}$}{$6\,800\,\text{g}$}{$68\,000\,\text{g}$}
\cauTN{10}{Năm $2026$ thuộc thế kỉ nào?}{Thế kỉ XIX}{Thế kỉ XX}{Thế kỉ XXI}{Thế kỉ XXII}
\cauTNhaiDong{11}{Một hình chữ nhật có chu vi $30\,\text{cm}$, chiều dài $9\,\text{cm}$. Diện tích là:}{$27\,\text{cm}^2$}{$54\,\text{cm}^2$}{$45\,\text{cm}^2$}{$36\,\text{cm}^2$}
\cauTNhaiDong{12}{Mua $4$ chiếc bánh ngọt hết $32\,000$ đồng. Mua $7$ chiếc bánh như thế hết số tiền là:}{$48\,000$ đồng}{$56\,000$ đồng}{$64\,000$ đồng}{$40\,000$ đồng}

\vspace{0.4mm}
\noindent\textbf{\large PHẦN II. TỰ LUẬN (7,0 điểm)}

\noindent\textbf{Bài 1 (2,0 điểm). Đặt tính rồi tính:}
\vspace{0.4mm}

\noindent
\begin{tabularx}{\linewidth}{|X|X|X|X|}
\hline
\oDatTinh[4]{a)}{$48\,356 + 36\,478$} & 
\oDatTinh[4]{b)}{$85\,310 - 47\,625$} & 
\oDatTinh[4]{c)}{$12\,345 \times 5$} & 
\oDatTinh[4]{d)}{$42\,816 : 6$} \tabularnewline
\hline
\end{tabularx}

\newpage
% -------------------- TRANG 2: MẶT SAU --------------------
\noindent
\begin{tabularx}{\linewidth}{|X|c|}
\hline
\textbf{PHẦN TỰ LUẬN (TIẾP THEO) -- PHIẾU BÀI TẬP TUẦN 09} & \textbf{MÃ ĐỀ: 404} \\
\hline
\end{tabularx}

\vspace{0.8mm}

% BÀI 2: TÌM X (4 Ý)
\khungBaiTap{%
  \textbf{Bài 2 (1,5 điểm). Tìm $x$:}\par\vspace{0.6mm}
  \noindent
  \begin{minipage}[t]{0.485\linewidth}
    \textbf{a)} $x - 19\,280 = 35\,640$\par
    \dongke{3}
    \vspace{0.4mm}
    \textbf{c)} $x \times 4 = 48\,936$\par
    \dongke{3}
  \end{minipage}\hfill
  \begin{minipage}[t]{0.485\linewidth}
    \textbf{b)} $82\,540 - x = 45\,120$\par
    \dongke{3}
    \vspace{0.4mm}
    \textbf{d)} $x : 9 = 5\,214$\par
    \dongke{3}
  \end{minipage}
}

\vspace{0.8mm}

% BÀI 3: ĐỔI ĐƠN VỊ (4 Ý - 1 CỘT TỪ TRÊN XUỐNG)
\khungBaiTap{%
  \textbf{Bài 3 (1,0 điểm). Viết số thích hợp vào chỗ chấm:}\par\vspace{0.8mm}
  \noindent
  a) $6\,\text{km}\,80\,\text{m} = \dotfill\,\text{m}$\par\vspace{1.4mm}
  \noindent
  b) $9\,\text{kg}\,350\,\text{g} = \dotfill\,\text{g}$\par\vspace{1.4mm}
  \noindent
  c) $5\,\text{m}^2\,72\,\text{dm}^2 = \dotfill\,\text{dm}^2$\par\vspace{1.4mm}
  \noindent
  d) $5\,\text{phút}\,20\,\text{giây} = \dotfill\,\text{giây}$\par
}

\vspace{0.8mm}

% BÀI 4: BÀI TOÁN GIẢI (11 DÒNG KẺ)
\khungBaiTap{%
  \textbf{Bài 4 (1,5 điểm).} Một thửa ruộng hình bình hành có độ dài đáy $250\,m$, chiều cao $120\,m$. Người ta cấy lúa trên thửa ruộng đó, cứ $100\,m^2$ thu hoạch được $70\,kg$ thóc. Hỏi thửa ruộng thu hoạch được bao nhiêu tạ thóc? Diện tích thửa ruộng bằng bao nhiêu héc-ta?
  \par\vspace{0.3mm}%
  \noindent\textit{Bài giải:}
  \dongke{11}
}

\vspace{0.8mm}

% BÀI 5: TOÁN NÂNG CAO (8 DÒNG KẺ)
\khungBaiTap{%
  \textbf{Bài 5 (1,0 điểm).} Tính bằng cách thuận tiện nhất:\par\vspace{0.6mm}\centerline{$K = 0{,}5 \times 12{,}5 \times 2 + 7{,}5 \times 10$}
  \par\vspace{0.3mm}%
  \noindent\textit{Bài làm:}
  \dongke{8}
}

\newpage


% =============================================================================
% ========================= MÃ ĐỀ: 405 ========================================
% =============================================================================

% -------------------- TRANG 1: MẶT TRƯỚC --------------------
\noindent
\begin{tabularx}{\linewidth}{|X|c|}
\hline
\textbf{PHIẾU BÀI TẬP CUỐI TUẦN -- TOÁN LỚP 4 (TUẦN 09)} & \textbf{MÃ ĐỀ: 405} \\
\textit{Trung tâm Toán Thầy Lĩnh -- Năm học: 2026 - 2027} & \textit{Thời gian: 45 phút} \\
\hline
\end{tabularx}

\vspace{1mm}

\noindent
\begin{tabularx}{\linewidth}{|X|p{3.5cm}|}
\hline
\rule[-2mm]{0pt}{7.5mm}\textbf{Họ và tên:} \dotfill & \textbf{Lớp:} \makebox[2cm]{\dotfill} \\
\hline
\end{tabularx}

\vspace{1mm}

\noindent
\begin{tabularx}{\linewidth}{|p{3cm}|X|}
\hline
\centering\textbf{ĐIỂM SỐ} & \centering\textbf{LỜI NHẬN XÉT CỦA THẦY CÔ} \tabularnewline
\hline
\rule{0pt}{1.1cm} & \tabularnewline
\hline
\end{tabularx}

\vspace{1mm}

\noindent\textbf{BẢNG ĐIỀN ĐÁP ÁN TRẮC NGHIỆM (Mỗi câu đúng 0,25 điểm):}
\vspace{0.6mm}

\begin{tabularx}{\linewidth}{|c|*{12}{>{\centering\arraybackslash}X|}}
\hline
\textbf{Câu} & \textbf{1} & \textbf{2} & \textbf{3} & \textbf{4} & \textbf{5} & \textbf{6} & \textbf{7} & \textbf{8} & \textbf{9} & \textbf{10} & \textbf{11} & \textbf{12} \\
\hline
\textbf{Đáp án} & & & & & & & & & & & & \\
\hline
\end{tabularx}

\vspace{1.2mm}

\noindent\textbf{\large PHẦN I. TRẮC NGHIỆM KHÁCH QUAN (3,0 điểm)}

\cauTN{1}{Số gồm $8$ chục nghìn, $3$ trăm và $5$ đơn vị viết là:}{$80\,305$}{$83\,005$}{$80\,350$}{$803\,005$}
\cauTN{2}{Chữ số $9$ trong số $59\,124$ thuộc hàng nào?}{Hàng chục}{Hàng trăm}{Hàng nghìn}{Hàng chục nghìn}
\cauTN{3}{Số liền trước của số $50\,000$ là:}{$49\,999$}{$50\,001$}{$49\,990$}{$49\,900$}
\cauTN{4}{Số lớn nhất trong các số $92\,645;\; 92\,654;\; 92\,456;\; 92\,564$ là:}{$92\,645$}{$92\,654$}{$92\,456$}{$92\,564$}
\cauTN{5}{Làm tròn số $58\,642$ đến hàng trăm ta được số:}{$58\,600$}{$58\,700$}{$59\,000$}{$58\,650$}
\cauTN{6}{Trong các số sau, số nào là số chẵn?}{$43\,571$}{$71\,432$}{$65\,809$}{$89\,215$}
\cauTN{7}{Giá trị của biểu thức $a + b$ với $a = 45\,000$ và $b = 25\,000$ là:}{$60\,000$}{$70\,000$}{$80\,000$}{$65\,000$}
\cauTN{8}{Số thích hợp điền vào chỗ chấm $7\,\text{m}\,3\,\text{cm} = \dots\,\text{cm}$ là:}{$73$}{$703$}{$730$}{$7\,003$}
\cauTN{9}{Khối lượng $3\,\text{kg}\,750\,\text{g} = \dots\,\text{g}$ là:}{$375\,\text{g}$}{$3\,075\,\text{g}$}{$3\,750\,\text{g}$}{$37\,500\,\text{g}$}
\cauTN{10}{Năm $1990$ thuộc thế kỉ nào?}{Thế kỉ XVIII}{Thế kỉ XIX}{Thế kỉ XX}{Thế kỉ XXI}
\cauTNhaiDong{11}{Một hình vuông có chu vi là $40\,\text{cm}$. Diện tích của hình vuông đó là:}{$80\,\text{cm}^2$}{$100\,\text{cm}^2$}{$40\,\text{cm}^2$}{$20\,\text{cm}^2$}
\cauTNhaiDong{12}{Mua $5$ cái thước kẻ hết $25\,000$ đồng. Mua $9$ cái thước kẻ như thế hết số tiền là:}{$35\,000$ đồng}{$45\,000$ đồng}{$50\,000$ đồng}{$40\,000$ đồng}

\vspace{0.4mm}
\noindent\textbf{\large PHẦN II. TỰ LUẬN (7,0 điểm)}

\noindent\textbf{Bài 1 (2,0 điểm). Đặt tính rồi tính:}
\vspace{0.4mm}

\noindent
\begin{tabularx}{\linewidth}{|X|X|X|X|}
\hline
\oDatTinh[4]{a)}{$37\,482 + 45\,239$} & 
\oDatTinh[4]{b)}{$73\,520 - 39\,465$} & 
\oDatTinh[4]{c)}{$14\,180 \times 6$} & 
\oDatTinh[4]{d)}{$49\,357 : 7$} \tabularnewline
\hline
\end{tabularx}

\newpage
% -------------------- TRANG 2: MẶT SAU --------------------
\noindent
\begin{tabularx}{\linewidth}{|X|c|}
\hline
\textbf{PHẦN TỰ LUẬN (TIẾP THEO) -- PHIẾU BÀI TẬP TUẦN 09} & \textbf{MÃ ĐỀ: 405} \\
\hline
\end{tabularx}

\vspace{0.8mm}

% BÀI 2: TÌM X (4 Ý)
\khungBaiTap{%
  \textbf{Bài 2 (1,5 điểm). Tìm $x$:}\par\vspace{0.6mm}
  \noindent
  \begin{minipage}[t]{0.485\linewidth}
    \textbf{a)} $x - 16\,540 = 28\,390$\par
    \dongke{3}
    \vspace{0.4mm}
    \textbf{c)} $x \times 6 = 54\,816$\par
    \dongke{3}
  \end{minipage}\hfill
  \begin{minipage}[t]{0.485\linewidth}
    \textbf{b)} $69\,480 - x = 34\,250$\par
    \dongke{3}
    \vspace{0.4mm}
    \textbf{d)} $x : 5 = 9\,342$\par
    \dongke{3}
  \end{minipage}
}

\vspace{0.8mm}

% BÀI 3: ĐỔI ĐƠN VỊ (4 Ý - 1 CỘT TỪ TRÊN XUỐNG)
\khungBaiTap{%
  \textbf{Bài 3 (1,0 điểm). Viết số thích hợp vào chỗ chấm:}\par\vspace{0.8mm}
  \noindent
  a) $4\,\text{km}\,35\,\text{m} = \dotfill\,\text{m}$\par\vspace{1.4mm}
  \noindent
  b) $7\,\text{kg}\,180\,\text{g} = \dotfill\,\text{g}$\par\vspace{1.4mm}
  \noindent
  c) $9\,\text{m}^2\,25\,\text{dm}^2 = \dotfill\,\text{dm}^2$\par\vspace{1.4mm}
  \noindent
  d) $4\,\text{phút}\,50\,\text{giây} = \dotfill\,\text{giây}$\par
}

\vspace{0.8mm}

% BÀI 4: BÀI TOÁN GIẢI (11 DÒNG KẺ)
\khungBaiTap{%
  \textbf{Bài 4 (1,5 điểm).} Một mảnh đất hình chữ nhật có chiều dài $180\,m$, chiều rộng bằng $\frac{5}{6}$ chiều dài. Người ta trồng ngô trên mảnh đất đó, trung bình mỗi héc-ta thu được $9$ tấn ngô. Hỏi cả mảnh đất thu hoạch được bao nhiêu tấn ngô? Diện tích mảnh đất bằng bao nhiêu héc-ta?
  \par\vspace{0.3mm}%
  \noindent\textit{Bài giải:}
  \dongke{11}
}

\vspace{0.8mm}

% BÀI 5: TOÁN NÂNG CAO (8 DÒNG KẺ)
\khungBaiTap{%
  \textbf{Bài 5 (1,0 điểm).} Tính bằng cách thuận tiện nhất:\par\vspace{0.6mm}\centerline{$K = 3{,}5 \times 8 - 8 \times 1{,}5$}
  \par\vspace{0.3mm}%
  \noindent\textit{Bài làm:}
  \dongke{8}
}

\newpage


% =============================================================================
% ========================= MÃ ĐỀ: 406 ========================================
% =============================================================================

% -------------------- TRANG 1: MẶT TRƯỚC --------------------
\noindent
\begin{tabularx}{\linewidth}{|X|c|}
\hline
\textbf{PHIẾU BÀI TẬP CUỐI TUẦN -- TOÁN LỚP 4 (TUẦN 09)} & \textbf{MÃ ĐỀ: 406} \\
\textit{Trung tâm Toán Thầy Lĩnh -- Năm học: 2026 - 2027} & \textit{Thời gian: 45 phút} \\
\hline
\end{tabularx}

\vspace{1mm}

\noindent
\begin{tabularx}{\linewidth}{|X|p{3.5cm}|}
\hline
\rule[-2mm]{0pt}{7.5mm}\textbf{Họ và tên:} \dotfill & \textbf{Lớp:} \makebox[2cm]{\dotfill} \\
\hline
\end{tabularx}

\vspace{1mm}

\noindent
\begin{tabularx}{\linewidth}{|p{3cm}|X|}
\hline
\centering\textbf{ĐIỂM SỐ} & \centering\textbf{LỜI NHẬN XÉT CỦA THẦY CÔ} \tabularnewline
\hline
\rule{0pt}{1.1cm} & \tabularnewline
\hline
\end{tabularx}

\vspace{1mm}

\noindent\textbf{BẢNG ĐIỀN ĐÁP ÁN TRẮC NGHIỆM (Mỗi câu đúng 0,25 điểm):}
\vspace{0.6mm}

\begin{tabularx}{\linewidth}{|c|*{12}{>{\centering\arraybackslash}X|}}
\hline
\textbf{Câu} & \textbf{1} & \textbf{2} & \textbf{3} & \textbf{4} & \textbf{5} & \textbf{6} & \textbf{7} & \textbf{8} & \textbf{9} & \textbf{10} & \textbf{11} & \textbf{12} \\
\hline
\textbf{Đáp án} & & & & & & & & & & & & \\
\hline
\end{tabularx}

\vspace{1.2mm}

\noindent\textbf{\large PHẦN I. TRẮC NGHIỆM KHÁCH QUAN (3,0 điểm)}

\cauTN{1}{Số gồm $4$ chục nghìn, $9$ trăm và $7$ đơn vị viết là:}{$40\,907$}{$49\,007$}{$40\,970$}{$409\,007$}
\cauTN{2}{Chữ số $6$ trong số $76\,215$ thuộc hàng nào?}{Hàng chục}{Hàng trăm}{Hàng nghìn}{Hàng chục nghìn}
\cauTN{3}{Số liền sau của số $69\,999$ là:}{$70\,000$}{$69\,998$}{$70\,001$}{$69\,900$}
\cauTN{4}{Số bé nhất trong các số $54\,321;\; 54\,123;\; 54\,213;\; 54\,312$ là:}{$54\,321$}{$54\,123$}{$54\,213$}{$54\,312$}
\cauTN{5}{Làm tròn số $83\,472$ đến hàng trăm ta được số:}{$83\,400$}{$83\,500$}{$83\,000$}{$83\,470$}
\cauTN{6}{Trong các số sau, số nào là số lẻ?}{$52\,840$}{$64\,126$}{$86\,531$}{$90\,708$}
\cauTN{7}{Giá trị của biểu thức $p - q$ với $p = 75\,000$ và $q = 30\,000$ là:}{$35\,000$}{$45\,000$}{$55\,000$}{$40\,000$}
\cauTN{8}{Số thích hợp điền vào chỗ chấm $9\,\text{m}\,4\,\text{cm} = \dots\,\text{cm}$ là:}{$94$}{$904$}{$940$}{$9\,004$}
\cauTN{9}{Khối lượng $8\,\text{kg}\,400\,\text{g} = \dots\,\text{g}$ là:}{$840\,\text{g}$}{$8\,040\,\text{g}$}{$8\,400\,\text{g}$}{$84\,000\,\text{g}$}
\cauTN{10}{Một giờ bằng bao nhiêu giây?}{$60$ giây}{$600$ giây}{$3\,600$ giây}{$360$ giây}
\cauTNhaiDong{11}{Một hình chữ nhật có chiều dài $15\,\text{cm}$, chiều rộng $6\,\text{cm}$. Diện tích là:}{$42\,\text{cm}^2$}{$90\,\text{cm}^2$}{$21\,\text{cm}^2$}{$60\,\text{cm}^2$}
\cauTNhaiDong{12}{Mua $6$ hộp chì màu hết $54\,000$ đồng. Mua $8$ hộp chì màu như thế hết số tiền là:}{$63\,000$ đồng}{$72\,000$ đồng}{$81\,000$ đồng}{$60\,000$ đồng}

\vspace{0.4mm}
\noindent\textbf{\large PHẦN II. TỰ LUẬN (7,0 điểm)}

\noindent\textbf{Bài 1 (2,0 điểm). Đặt tính rồi tính:}
\vspace{0.4mm}

\noindent
\begin{tabularx}{\linewidth}{|X|X|X|X|}
\hline
\oDatTinh[4]{a)}{$52\,648 + 38\,175$} & 
\oDatTinh[4]{b)}{$86\,240 - 48\,375$} & 
\oDatTinh[4]{c)}{$13\,416 \times 7$} & 
\oDatTinh[4]{d)}{$37\,548 : 4$} \tabularnewline
\hline
\end{tabularx}

\newpage
% -------------------- TRANG 2: MẶT SAU --------------------
\noindent
\begin{tabularx}{\linewidth}{|X|c|}
\hline
\textbf{PHẦN TỰ LUẬN (TIẾP THEO) -- PHIẾU BÀI TẬP TUẦN 09} & \textbf{MÃ ĐỀ: 406} \\
\hline
\end{tabularx}

\vspace{0.8mm}

% BÀI 2: TÌM X (4 Ý)
\khungBaiTap{%
  \textbf{Bài 2 (1,5 điểm). Tìm $x$:}\par\vspace{0.6mm}
  \noindent
  \begin{minipage}[t]{0.485\linewidth}
    \textbf{a)} $x - 27\,340 = 38\,520$\par
    \dongke{3}
    \vspace{0.4mm}
    \textbf{c)} $x \times 8 = 64\,832$\par
    \dongke{3}
  \end{minipage}\hfill
  \begin{minipage}[t]{0.485\linewidth}
    \textbf{b)} $91\,620 - x = 46\,380$\par
    \dongke{3}
    \vspace{0.4mm}
    \textbf{d)} $x : 6 = 8\,415$\par
    \dongke{3}
  \end{minipage}
}

\vspace{0.8mm}

% BÀI 3: ĐỔI ĐƠN VỊ (4 Ý - 1 CỘT TỪ TRÊN XUỐNG)
\khungBaiTap{%
  \textbf{Bài 3 (1,0 điểm). Viết số thích hợp vào chỗ chấm:}\par\vspace{0.8mm}
  \noindent
  a) $8\,\text{km}\,50\,\text{m} = \dotfill\,\text{m}$\par\vspace{1.4mm}
  \noindent
  b) $4\,\text{kg}\,920\,\text{g} = \dotfill\,\text{g}$\par\vspace{1.4mm}
  \noindent
  c) $4\,\text{m}^2\,68\,\text{dm}^2 = \dotfill\,\text{dm}^2$\par\vspace{1.4mm}
  \noindent
  d) $3\,\text{phút}\,25\,\text{giây} = \dotfill\,\text{giây}$\par
}

\vspace{0.8mm}

% BÀI 4: BÀI TOÁN GIẢI (11 DÒNG KẺ)
\khungBaiTap{%
  \textbf{Bài 4 (1,5 điểm).} Một thửa ruộng hình chữ nhật có chiều dài $160\,m$, chiều rộng $125\,m$. Người ta cấy lúa trên thửa ruộng đó, cứ $100\,m^2$ thu hoạch được $62\,kg$ thóc. Hỏi thửa ruộng thu hoạch được bao nhiêu tạ thóc? Diện tích thửa ruộng bằng bao nhiêu héc-ta?
  \par\vspace{0.3mm}%
  \noindent\textit{Bài giải:}
  \dongke{11}
}

\vspace{0.8mm}

% BÀI 5: TOÁN NÂNG CAO (8 DÒNG KẺ)
\khungBaiTap{%
  \textbf{Bài 5 (1,0 điểm).} Tính bằng cách thuận tiện nhất:\par\vspace{0.6mm}\centerline{$K = 0{,}4 \times 25 + 0{,}4 \times 75$}
  \par\vspace{0.3mm}%
  \noindent\textit{Bài làm:}
  \dongke{8}
}

\newpage


% =============================================================================
% ========================= MÃ ĐỀ: 407 ========================================
% =============================================================================

% -------------------- TRANG 1: MẶT TRƯỚC --------------------
\noindent
\begin{tabularx}{\linewidth}{|X|c|}
\hline
\textbf{PHIẾU BÀI TẬP CUỐI TUẦN -- TOÁN LỚP 4 (TUẦN 09)} & \textbf{MÃ ĐỀ: 407} \\
\textit{Trung tâm Toán Thầy Lĩnh -- Năm học: 2026 - 2027} & \textit{Thời gian: 45 phút} \\
\hline
\end{tabularx}

\vspace{1mm}

\noindent
\begin{tabularx}{\linewidth}{|X|p{3.5cm}|}
\hline
\rule[-2mm]{0pt}{7.5mm}\textbf{Họ và tên:} \dotfill & \textbf{Lớp:} \makebox[2cm]{\dotfill} \\
\hline
\end{tabularx}

\vspace{1mm}

\noindent
\begin{tabularx}{\linewidth}{|p{3cm}|X|}
\hline
\centering\textbf{ĐIỂM SỐ} & \centering\textbf{LỜI NHẬN XÉT CỦA THẦY CÔ} \tabularnewline
\hline
\rule{0pt}{1.1cm} & \tabularnewline
\hline
\end{tabularx}

\vspace{1mm}

\noindent\textbf{BẢNG ĐIỀN ĐÁP ÁN TRẮC NGHIỆM (Mỗi câu đúng 0,25 điểm):}
\vspace{0.6mm}

\begin{tabularx}{\linewidth}{|c|*{12}{>{\centering\arraybackslash}X|}}
\hline
\textbf{Câu} & \textbf{1} & \textbf{2} & \textbf{3} & \textbf{4} & \textbf{5} & \textbf{6} & \textbf{7} & \textbf{8} & \textbf{9} & \textbf{10} & \textbf{11} & \textbf{12} \\
\hline
\textbf{Đáp án} & & & & & & & & & & & & \\
\hline
\end{tabularx}

\vspace{1.2mm}

\noindent\textbf{\large PHẦN I. TRẮC NGHIỆM KHÁCH QUAN (3,0 điểm)}

\cauTN{1}{Số gồm $6$ chục nghìn, $7$ trăm và $4$ đơn vị viết là:}{$60\,704$}{$67\,004$}{$60\,740$}{$607\,004$}
\cauTN{2}{Chữ số $3$ trong số $83\,615$ thuộc hàng nào?}{Hàng trăm}{Hàng chục}{Hàng nghìn}{Hàng chục nghìn}
\cauTN{3}{Số liền trước của số $80\,000$ là:}{$79\,999$}{$80\,001$}{$79\,990$}{$79\,900$}
\cauTN{4}{Số lớn nhất trong các số $73\,845;\; 73\,854;\; 73\,485;\; 73\,584$ là:}{$73\,845$}{$73\,854$}{$73\,485$}{$73\,584$}
\cauTN{5}{Làm tròn số $67\,824$ đến hàng trăm ta được số:}{$67\,800$}{$67\,900$}{$68\,000$}{$67\,820$}
\cauTN{6}{Trong các số sau, số nào là số chẵn?}{$31\,759$}{$92\,568$}{$45\,301$}{$67\,823$}
\cauTN{7}{Giá trị của biểu thức $u + v$ với $u = 52\,000$ và $v = 38\,000$ là:}{$80\,000$}{$90\,000$}{$100\,000$}{$85\,000$}
\cauTN{8}{Số thích hợp điền vào chỗ chấm $3\,\text{m}\,7\,\text{cm} = \dots\,\text{cm}$ là:}{$37$}{$307$}{$370$}{$3\,007$}
\cauTN{9}{Khối lượng $9\,\text{kg}\,150\,\text{g} = \dots\,\text{g}$ là:}{$915\,\text{g}$}{$9\,015\,\text{g}$}{$9\,150\,\text{g}$}{$91\,500\,\text{g}$}
\cauTN{10}{Thế kỉ XX có bao nhiêu năm?}{$10$ năm}{$50$ năm}{$100$ năm}{$1\,000$ năm}
\cauTNhaiDong{11}{Một hình vuông có chu vi là $28\,\text{cm}$. Diện tích hình vuông đó là:}{$56\,\text{cm}^2$}{$49\,\text{cm}^2$}{$28\,\text{cm}^2$}{$14\,\text{cm}^2$}
\cauTNhaiDong{12}{Mua $7$ chai nước khoáng hết $42\,000$ đồng. Mua $9$ chai nước khoáng như thế hết số tiền là:}{$48\,000$ đồng}{$54\,000$ đồng}{$63\,000$ đồng}{$60\,000$ đồng}

\vspace{0.4mm}
\noindent\textbf{\large PHẦN II. TỰ LUẬN (7,0 điểm)}

\noindent\textbf{Bài 1 (2,0 điểm). Đặt tính rồi tính:}
\vspace{0.4mm}

\noindent
\begin{tabularx}{\linewidth}{|X|X|X|X|}
\hline
\oDatTinh[4]{a)}{$46\,835 + 37\,549$} & 
\oDatTinh[4]{b)}{$94\,310 - 58\,475$} & 
\oDatTinh[4]{c)}{$15\,214 \times 5$} & 
\oDatTinh[4]{d)}{$65\,439 : 9$} \tabularnewline
\hline
\end{tabularx}

\newpage
% -------------------- TRANG 2: MẶT SAU --------------------
\noindent
\begin{tabularx}{\linewidth}{|X|c|}
\hline
\textbf{PHẦN TỰ LUẬN (TIẾP THEO) -- PHIẾU BÀI TẬP TUẦN 09} & \textbf{MÃ ĐỀ: 407} \\
\hline
\end{tabularx}

\vspace{0.8mm}

% BÀI 2: TÌM X (4 Ý)
\khungBaiTap{%
  \textbf{Bài 2 (1,5 điểm). Tìm $x$:}\par\vspace{0.6mm}
  \noindent
  \begin{minipage}[t]{0.485\linewidth}
    \textbf{a)} $x - 24\,180 = 43\,560$\par
    \dongke{3}
    \vspace{0.4mm}
    \textbf{c)} $x \times 7 = 49\,644$\par
    \dongke{3}
  \end{minipage}\hfill
  \begin{minipage}[t]{0.485\linewidth}
    \textbf{b)} $85\,430 - x = 37\,210$\par
    \dongke{3}
    \vspace{0.4mm}
    \textbf{d)} $x : 7 = 9\,216$\par
    \dongke{3}
  \end{minipage}
}

\vspace{0.8mm}

% BÀI 3: ĐỔI ĐƠN VỊ (4 Ý - 1 CỘT TỪ TRÊN XUỐNG)
\khungBaiTap{%
  \textbf{Bài 3 (1,0 điểm). Viết số thích hợp vào chỗ chấm:}\par\vspace{0.8mm}
  \noindent
  a) $3\,\text{km}\,95\,\text{m} = \dotfill\,\text{m}$\par\vspace{1.4mm}
  \noindent
  b) $6\,\text{kg}\,650\,\text{g} = \dotfill\,\text{g}$\par\vspace{1.4mm}
  \noindent
  c) $8\,\text{m}^2\,42\,\text{dm}^2 = \dotfill\,\text{dm}^2$\par\vspace{1.4mm}
  \noindent
  d) $5\,\text{phút}\,45\,\text{giây} = \dotfill\,\text{giây}$\par
}

\vspace{0.8mm}

% BÀI 4: BÀI TOÁN GIẢI (11 DÒNG KẺ)
\khungBaiTap{%
  \textbf{Bài 4 (1,5 điểm).} Một mảnh đất hình bình hành có độ dài đáy $200\,m$, chiều cao $120\,m$. Người ta trồng lạc trên mảnh đất đó, trung bình mỗi héc-ta thu được $4{,}5$ tấn lạc. Hỏi cả mảnh đất thu hoạch được bao nhiêu tấn lạc? Diện tích mảnh đất bằng bao nhiêu héc-ta?
  \par\vspace{0.3mm}%
  \noindent\textit{Bài giải:}
  \dongke{11}
}

\vspace{0.8mm}

% BÀI 5: TOÁN NÂNG CAO (8 DÒNG KẺ)
\khungBaiTap{%
  \textbf{Bài 5 (1,0 điểm).} Tính bằng cách thuận tiện nhất:\par\vspace{0.6mm}\centerline{$K = 1{,}25 \times 3{,}2 + 1{,}25 \times 6{,}8 - 2{,}5$}
  \par\vspace{0.3mm}%
  \noindent\textit{Bài làm:}
  \dongke{8}
}

\newpage
\end{document}
